% ===== §16 =====
\section{The Dialectic of Cultural Evolution}
\label{sec:evolution}

\noindent
One question remains before the synthesis is complete and the paper can descend into the intimate dyad: how culture, so defined, \emph{changes}. The definition gave culture's mode of being and the sublation of power secured its justice, but a culture that is generation rather than deposit must have an account of its own movement through time, and that account must avoid two opposite errors as carefully as the dialectic of the material and the ideal avoided its two. This section gives the dialectic of cultural evolution, and in giving it completes both the synthesis and the Hegelian line that has run through the whole part.

\subsection*{Two errors about cultural change}

The first error is the model of \emph{linear progress}: the idea, inherited from the evolutionary anthropology that stands behind the very first tradition, that culture develops along a single ascending track from lower to higher forms, each stage an advance on the last, the whole tending toward a terminus of completion or truth. This is the error that reads the history of culture as a march, and it is refused by everything the paper has established: if the truth of culture is historical and non-terminal, if every cultural form is one materially conditioned capture of a Real that overflows it, then there is no single track and no terminus, and the notion of a highest culture toward which all others are steps is a projection of one form's self-image onto the whole of history. The second error is the opposite: the model of \emph{mere reproduction}, the pessimistic determinism that reads cultural change as no change at all, only the endless recirculation of a power order that reproduces itself under cosmetic variation, so that what looks like development is always the same domination in a new dress. This is the error to which the power tradition, taken without its dialectical corrective, tends: if culture is nothing but the encoding of power, then its history is nothing but the reproduction of power, and the appearance of the new is an illusion. The first error sees only ascent; the second sees only the wheel. Both miss the actual movement, which is neither a line nor a circle but a dialectic.

\subsection*{Evolution as the contradiction of power-generation and counter-power-generation}

The actual movement can be stated as a contradiction, and it is the contradiction the previous section prepared. Cultural evolution is the movement of \emph{power-generation against counter-power-generation}. Every cultural form, in generating value, generates power, and that power tends toward entrenchment---toward congealing into a fixed order, a hierarchy naturalized, a hegemony settled. This is the moment of truth in the reproduction model: left to itself, the generation of culture would indeed tend toward the reproduction of an entrenching power. But the generation of power is met, in a just generation, by the internal generation of counter-power, the friction and the re-orientation of the previous section; and this counter-power is the negation that the entrenching power calls forth from within the same process. The culture generates the power; the power tends to entrench; the generation, if just, generates against its own entrenching power a counter-power that negates the entrenchment; and this negation, in turn, opens the culture to a new form, which generates its own power, which tends to entrench, which calls forth its own counter-power, and so on. This is not a line, because it does not ascend toward a terminus; it is not a circle, because the counter-power does not return the culture to where it began but opens it to a form that did not exist before. It is a spiral of determinate negation, and its motor, as always in this paper, is material: which powers entrench and which counter-powers can be generated against them are set by the forces and relations of production, so that the dialectic of cultural evolution is a materially driven contradiction and not a self-movement of spirit.

\subsection*{The negation of the negation, materialized}

This is the point at which the Hegelian line of the paper reaches its completion, and it should be named as such. The movement just described is the negation of the negation, the deepest figure of the dialectic, here given a material body and a cultural content. The first negation is the entrenchment of power: the generated culture goes out of itself into a power that congeals, that stands over against the generativity that produced it, that threatens to fix the living relation into a dead hierarchy---the moment of estrangement, of exteriorization become domination. The second negation, the negation of this negation, is the counter-power generated from within: not a return to an original innocence before power, which never existed, but a determinate negation of the entrenchment that carries the culture forward to a new form, a form that preserves what the entrenched power had generated while negating its fixity. Estrangement is thus not the end of the story, as the reproduction model would have it, but a moment in a movement that passes through it; the exteriorized culture that has become a dominating power can be negated in its turn and taken up into a form in which the relation reclaims what its own power had captured. This is why the paper can hold, at once, that power is internal to all culture and that culture is not doomed to the reproduction of domination: because the internality of power is the first negation, and the internal generation of counter-power is the second, and the second is as native to a just generation as the first is to generation as such.

\subsection*{Evaluating a culture historically}

From this follows the mode of evaluation the fourth revision announced, and stating it completes the synthesis. A cultural form is to be evaluated neither by its distance along a track of progress nor by its mere function in reproducing power, but dialectically, by its place in this movement. One asks of a form both what it generated and how its power stood: whether, under the material conditions that produced it, it was a genuine advance, a capture of the Real that let the relation generate more than it could before---and whether the power it generated was met by a counter-power or allowed to entrench. The same form may be, under its own conditions, a real progress, and, when conditions change and the Real it captured overflows anew, a fetter to be negated. To criticize an old cultural form and to acknowledge its former progressiveness are therefore, as the programmatic claim held, two faces of one dialectical evaluation: the form was the negation its moment required, and it becomes, in a later moment, the negation to be negated.

From this the paper draws a fourth thesis of its own, a criterion of cultural persistence that stands alongside the geometric criterion of the good cycle developed earlier in the series. Call it the \emph{counter-power criterion}: whether a cultural form can continue to generate across a change in the conditions that produced it turns on whether it generates, endogenously, a counter-power against the power it itself carries. A form that transmits its power without transmitting the means of that power's revision entrenches, and when its conditions change it passes from a form that generated into a form that fetters; a form that carries, with its power, a counter-power that keeps that power open to re-orientation can be turned to conditions its authors did not foresee, and so continues to generate where the rigid form breaks. The criterion does not equate persistence with the presence of counter-power, as though the one were the other; it holds that the endogenous generation of counter-power is what conditions the possibility of persistence through change, the difference between a form that bends and a form that breaks. It is the temporal and evolutionary counterpart, at the scale of a culture's history, of the requirement the sublation of power stated for a culture's justice: the just culture generates the counter-power that keeps its power from entrenching, and the persisting culture is the one whose counter-power lets it be re-oriented rather than broken when its moment passes.

This is the disciplined historical eye the paper has claimed throughout, neither the celebration of a march nor the cynicism of the wheel, but the evaluation of each form by its place in a materially driven spiral that has direction without a terminus. With this, the synthesis is complete: culture is defined, its power is answered, and its evolution is given its dialectic. The general ontology and epistemology of cultural emergence are established, and the paper turns to the particular for which the general was built---the emergence of culture in the intimate dyad.
